% =====================================================================
\chapter{Simuladores diferenciables y modelos sustitutos}
\label{tema:3}
% =====================================================================

El \cref{tema:2} obtuvo el gradiente de un simulador de elementos finitos derivando a mano sus ecuaciones. Este tema generaliza la idea a cualquier simulador, estacionario o temporal, y presenta la alternativa cuando derivar no es posible o el simulador es demasiado caro: aprender un modelo sustituto con pocas simulaciones y optimizar sobre él. Se siguen las notas de S.~G. Johnson sobre adjuntos para recurrencias~\cite{johnson}, las de cálculo matricial del MIT~\cite[§6.3]{mit18s096}, los capítulos de modelos sustitutos de Kochenderfer y Wheeler~\cite{kw}, y los artículos de DiffTaichi~\cite{difftaichi}, las redes informadas por la física~\cite{pinn}, el operador neuronal de Fourier~\cite{fno} y MeshGraphNets~\cite{mgn}. La física de la aleta procede del capítulo~3 de Incropera \emph{et al.}~\cite{incropera}.

\section{Derivar a través de un simulador}

Casi todo simulador discreto resuelve un sistema de ecuaciones $R(u,\theta)=0$, donde $u\in\R^m$ es el estado (temperaturas, desplazamientos, velocidades) y $\theta\in\R^n$ son los parámetros de diseño.

\begin{teorema}[Sensibilidad por la función implícita y el adjunto]\label{teo:implicita}
Sean $R:\R^m\times\R^n\to\R^m$ y $J(\theta)=g\bigl(u(\theta),\theta\bigr)$ funciones continuamente derivables, con $R(u(\theta),\theta)=0$ y $\partial R/\partial u$ invertible en la solución. Entonces
\[
\frac{du}{d\theta}=-\Bigl(\frac{\partial R}{\partial u}\Bigr)^{-1}\frac{\partial R}{\partial\theta},\qquad
\frac{dJ}{d\theta}=\frac{\partial g}{\partial\theta}-\lambda^\T\frac{\partial R}{\partial\theta},\quad\text{con}\quad\Bigl(\frac{\partial R}{\partial u}\Bigr)^{\!\T}\lambda=\Bigl(\frac{\partial g}{\partial u}\Bigr)^{\!\T}.
\]
\end{teorema}

El \cref{teo:adjunto} es el caso $R=K(x)U-F$. El coste del adjunto es una resolución con la matriz transpuesta, que en problemas lineales puede reutilizar la factorización ya calculada; el de la sensibilidad directa, $n$ resoluciones. Para simuladores temporales, $u_{k+1}=\Phi(u_k,\theta)$, el mismo razonamiento da una recurrencia adjunta que se recorre hacia atrás en el tiempo~\cite{johnson}:
\[
\lambda_N=\Bigl(\frac{\partial g}{\partial u_N}\Bigr)^{\!\T},\qquad \lambda_{k-1}=\Bigl(\frac{\partial\Phi}{\partial u}\Bigr)^{\!\T}_{\!k}\lambda_k,\qquad \frac{dg}{d\theta}=\frac{\partial g}{\partial\theta}+\sum_{k=1}^N\lambda_k^\T\frac{\partial\Phi}{\partial\theta}\Big|_{u_{k-1}} .
\]
Es la retropropagación a través del tiempo de una red recurrente (módulo~12, §5.3), con los mismos problemas: hay que guardar todos los estados, lo que se alivia con el recálculo por puntos de control (módulo~10, §5.2), y en dinámicas caóticas o con muchos pasos el gradiente puede explotar.

\section{Diferenciación automática de simuladores}

En lugar de derivar a mano, el simulador se puede escribir en un lenguaje que propague derivadas. JAX transforma programas de Python y NumPy (derivación, vectorización y compilación para GPU y TPU)~\cite{jax}; NVIDIA Warp compila núcleos de Python para GPU que son derivables~\cite{warp}; Taichi, con su extensión DiffTaichi, permite escribir simuladores físicos derivables, y sus autores informan de que su simulador elástico necesita unas 4,2 veces menos líneas que una versión CUDA escrita a mano, con velocidad similar, y es 188 veces más rápido que la implementación en TensorFlow~\cite{difftaichi}. PyTorch también sirve cuando el simulador se expresa con operaciones de tensores, como en el caso de estudio de este tema.

Tres advertencias limitan el enfoque. Los eventos discretos (contactos, cambios de malla, condicionales sobre el estado) producen funciones discontinuas cuya derivada automática es correcta pero inútil, porque vale cero o no ve el salto. Las derivadas de un resolvedor iterativo no deben calcularse a través de sus iteraciones, sino con el \cref{teo:implicita} en la solución convergida. Y el coste de memoria del modo inverso crece con el número de pasos de tiempo.

\section{Modelos sustitutos}

\begin{definicion}[Modelo sustituto]\label{def:sustituto}
Sea $s:\Theta\to\R$ una respuesta cara de evaluar (el resultado de un simulador). Un \emph{modelo sustituto} es una función $\hat s_\varphi$ barata, ajustada con un conjunto de diseño $\{(\theta_i,s(\theta_i))\}_{i=1}^M$, con la que se sustituye a $s$ dentro de la optimización.
\end{definicion}

Los sustitutos clásicos son las superficies de respuesta polinómicas ajustadas por mínimos cuadrados y los procesos gaussianos, que además estiman su incertidumbre~\cite{kw}. Los neuronales cubren desde un perceptrón que aproxima una respuesta escalar hasta modelos que aproximan campos completos: las redes informadas por la física (PINN) minimizan el residuo de la ecuación diferencial en puntos de colocación, calculando las derivadas de la red por diferenciación automática~\cite{pinn}; los operadores neuronales de Fourier aprenden la aplicación entre funciones de entrada y de salida, y sus autores informan de aceleraciones de hasta tres órdenes de magnitud frente a resolvedores clásicos~\cite{fno}; MeshGraphNets aprende la dinámica sobre la malla de simulación con una red de grafos~\cite{mgn}. Bibliotecas como NVIDIA PhysicsNeMo (antes Modulus) agrupan estas arquitecturas para construir sustitutos de CFD, mecánica estructural o análisis térmico~\cite{physicsnemo}.

El diseño de experimentos importa tanto como el modelo. El muestreo por hipercubo latino divide el rango de cada variable en $M$ intervalos y toma exactamente una muestra por intervalo, de modo que sus proyecciones sobre cada eje quedan uniformemente repartidas~\cite{kw}.

\begin{proposicion}[Error del óptimo de un sustituto]\label{prop:2eps}
Sea $\Theta$ compacto, $s$ y $\hat s$ continuas y $\sup_{\theta\in\Theta}\abs{\hat s(\theta)-s(\theta)}\le\varepsilon$. Si $\theta^\star$ minimiza $s$ y $\hat\theta$ minimiza $\hat s$ en $\Theta$, entonces
\[
s(\theta^\star)\le s(\hat\theta)\le s(\theta^\star)+2\varepsilon .
\]
\end{proposicion}

La proposición es tranquilizadora y peligrosa a la vez: la garantía exige un error \emph{uniforme} en todo $\Theta$, mientras que la validación solo mide el error en unos puntos. El optimizador busca precisamente las zonas en que $\hat s$ es optimista, que suelen ser aquellas en que hay pocos datos. Es el mismo fenómeno que la sobreoptimización de un modelo de recompensa (módulo~21, §6.4). Por eso todo diseño propuesto por un sustituto se verifica con el simulador y, si hay presupuesto, se añade a los datos y se reajusta el modelo.

\section{Caso de estudio: el perfil de una aleta disipadora}

\addtocounter{definicion}{1}%

\begin{table}[h]
\centering\small
\begin{tabularx}{\textwidth}{@{}l>{\raggedright\arraybackslash}X>{\raggedright\arraybackslash}p{2.3cm}>{\raggedright\arraybackslash}X@{}}
\toprule
Herramienta & Para qué sirve & Licencia & Cuándo elegirla \\
\midrule
JAX & derivación, vectorización y compilación de código NumPy & Apache-2.0 & simuladores nuevos escritos en estilo funcional \\
NVIDIA Warp & núcleos de simulación derivables en GPU & Apache-2.0 & física de partículas, telas o sólidos en GPU \\
Taichi (DiffTaichi) & lenguaje para simulación física derivable & Apache-2.0 & simuladores de partículas y rejillas de alto rendimiento \\
PyTorch & tensores y diferenciación automática & BSD-3 (con componentes de otras licencias) & simuladores pequeños unidos a redes neuronales \\
NVIDIA PhysicsNeMo & PINN, operadores neuronales y GNN para física & Apache-2.0 & sustitutos de CFD o de estructuras a escala industrial \\
scikit-learn & MLP, procesos gaussianos, validación & BSD-3 & sustitutos de pocas variables con pocos datos \\
dolfin-adjoint & adjuntos automáticos de modelos FEniCS & LGPL & EDP generales con elementos finitos \\
\bottomrule
\end{tabularx}
\caption{Tecnologías del tema~3. Licencias según los metadatos publicados por cada proyecto.}
\label{tab:tec3}
\end{table}

\addtocounter{definicion}{1}%

